\section{What remains hypothetical and how to proceed}

The results established here resolve specific questions for specific
models and classes of derivations. The following table records questions
that must not be listed as solved.

\begin{longtable}{L{0.30\textwidth}L{0.64\textwidth}}
\toprule Problem & Precise status and next verifiable step \\ \midrule
\endhead
All harmonic embeddings & Two infinite families and finite checks
are provided. A general passage between residue classes or a new
constructive rectification is needed. A bounded search is not a classification.\\[6pt]
Complex rectification of $H(3,7,8)$ and $H(4,5,7)$ & Not constructed.
One can search for a locally finite field in the affine spaces of
candidates obtained here. Absence of such a field at these degrees, if proved,
would not exclude fields of higher degree.\\[6pt]
Full ML of the new kernels & At present there is a dichotomy: $k$ or $k[t]$.
For the value $k$, one LND moving $t$ suffices.
For the value $k[t]$, one must prove that every LND fixes $t$;
a bounded search cannot prove this.\\[6pt]
Homogeneous $E$ with slice $t$ & The conditions $E(t)=1$ and preservation
of the kernel do not yet imply local nilpotence. Vanishing of iterates
on generators is an exact positive certificate; failure of finitely many
iterates to vanish does not prove that no suitable field exists.\\[6pt]
Formal coordinates & All finite thickenings and the completion have been
constructed. Termination of the coordinates as polynomials has not been
proved: this is a separate global problem.\\[6pt]
Dasgupta--Gaifullin Question 4.21 & Filtration formulas have been obtained,
but minimality of a three-step birational presentation has not
been established. All admissible shorter presentations must be excluded.\\[6pt]
Minimum degrees of rectifiers & The degrees of the exhibited automorphisms
are upper bounds. Minimum slice degrees do not replace this problem.\\[6pt]
Real descent & Real rectification of the nontrivial knots under consideration
is impossible. This alone does not exclude polynomiality of the corresponding
real kernel without preservation of the distinguished $t$.\\[6pt]
General cancellation and linearization problems & One proposed
candidate has been ruled out; the general problems and all possible
actions have not been fully investigated.\\
\bottomrule
\end{longtable}

\subsection*{Why the search for an LND with a slice is motivated}
In $B_\mu$, the element $t$ is distinguished by the embedding itself, while
the special fiber $t=0$ already has a polynomial description. If one can exhibit
a homogeneous derivation $E$ of degree $-1$ with $E(t)=1$ and prove
its local nilpotence, the slice formula gives
\[
 \rho_E(f)=\sum_{j\ge0}\frac{(-t)^j}{j!}E^j(f),\qquad
 B_\mu=(\ker E)[t].
\]
The sum is finite precisely because of local nilpotence. Under
homogeneity, this construction recovers coordinates for the original embedding,
as proved in \S11. Thus the additional LND requirement serves a precise
geometric purpose. It is not a technical detail that follows automatically
from the linear equations on coefficients.

For a finite generating set $b_1,\ldots,b_m$, it suffices to exhibit
integers $N_i$ and identities $E^{N_i}(b_i)=0$. The Leibniz rule then gives
local nilpotence on the whole algebra. For the search for a normal
locally finite field, an alternative positive certificate is a
finite-dimensional stable space containing the generators, as in \S12.
Both methods yield a verifiable proof when successful,
but no general time or degree bound for these searches is proved here.

\subsection*{The role of Cremona groups}
Birational changes of the form
\[
 (x,y,z)\longmapsto (x,y,m(x)z-q(x,y))
\]
provide an effective language for transferring fields and coordinates after
localization. Returning to the polynomial ring requires exact divisibility
and a check of the modification center. The obstruction established for
the natural center of the family $H(3,3n+1,3n+2)$ shows that merely having
a rational inverse is insufficient. The practical programme consists of
searching for other centers, computing image ideals and comparing
different factorizations. No general classification of Cremona subgroups
is claimed here.

\subsection*{Connection with the Abhyankar--Moh theorem}
A direct application of the planar theorem becomes possible once
a coordinate surface containing the curve has been exhibited: on this
surface, the embedded line becomes a planar embedding to which
\cite{AM,Suzuki} applies. Finding such a surface in the ambient space
is the substantive missing step. The equivalent criteria through
a coordinate parameter, a graded LND and a normal locally finite
field offer different methods for addressing precisely this step.
It must not be replaced by an appeal to the already known planar theorem.

\subsection*{Suggested reading sequence}
For an introduction to LNDs, a useful starting point is \emph{First Principles}
and \emph{Further Properties of LNDs} in Freudenburg's book \cite{Freudenburg},
followed by the chapters \emph{Dimension Three}, \emph{Algorithms},
\emph{Makar-Limanov and Derksen Invariants} and
\emph{Slices, Embeddings and Cancellation}. For this project,
the next readings are \cite{vdER} (the precise conjectures and the kernel ring),
\cite{Masuda} (equivariant structure), \cite{EVM} (relative
coordinates), \cite{DG} (iterations of LNDs and complexity),
and \cite{FF,KP} (harmonic curves and knots).

\section{Reproducibility and contents of the supplement}

All computational files are in \file{computations/}.
Comments in the supplied code are in English. The original
versions are retained in the provenance archive; a separate check verifies
that edits to Python comments and documentation strings do not change
the abstract syntax tree of the executable part.
This is a useful packaging check, but not a proof of algorithmic correctness.

\begin{longtable}{L{0.29\textwidth}L{0.65\textwidth}}
\toprule Directory & Purpose \\ \midrule
\endhead
\file{baseline/} & Original rectifiers; exact checks of mutually inverse
maps, printed components and original certificates.\\[5pt]
\file{legacy_appendix/} & Computations from the published appendix
of 26 October, including local dependencies for offline execution.\\[5pt]
\file{continuation/} & Fourteen continuation checks: minimum slice
degrees, multiplicity-dependent degrees, formal coordinates, transfer
restrictions, LND filtrations and relative ML.\\[5pt]
\file{euler/} & Normal fields, ranks of linear systems, dual
inconsistency certificates and invariant-axis obstructions for selected fields.\\[5pt]
\file{audit_checks/} & Verification identifying the candidate in
Conjectures 5.1--5.3 with the original model and original slice.\\[5pt]
\file{verification_logs/} & Standard output, errors and execution time
for every program in the repeated run.\\[5pt]
\file{provenance/} & Original versions, checksums and information
on the origin of the materials.\\
\bottomrule
\end{longtable}

The entry point is \file{computations/run_all_checks.py}.
Run the following from the \file{computations/} directory:
\begin{quote}\ttfamily\small
python -m pip install -r requirements.txt\\
python run\_all\_checks.py
\end{quote}
The fresh run of 28 October 2026 completed successfully: \textbf{35 out of 35
programs}. All principal algebraic checks are exact, without replacing
identities by numerical approximations. The final machine-readable record
\file{verification_summary.json} records the results and commands.
The number of programs is not the number of theorems: one program may
check many identities, while infinite families require the mathematical
proofs given above.

GAP/Macaulay2 code and historical logs are included for independent
checking, but those systems were not run in the current environment.
Their older successful logs are explicitly distinguished from the new Python run.
Lengthy search experiments are retained as source code and data;
they do not replace the shorter verification programs.

\subsection*{Compatibility with the repository}
The baseline was checked at commit
\file{b954f3609e900761fddc9c8ff03d8301e9604cab}.
All 75 files of the published appendix matched the Git objects of that
version. The update adds only a separate directory,
\file{updates/2026-10-28/}, and a pointer to it within the earlier
appendix. The main preprint and earlier certificates are not replaced.

Two earlier research directions are clarified substantively.
First, the relative ML invariant in the earlier appendix was proposed
conditionally; \S10 shows that the anticipated nontriviality is impossible
under that definition. Second, the new degree bound $\ge5$ for $H(4,5,7)$
concerns coordinates preserving one of the original variables.
The general search for an absolute coordinate starting at degree $4$ remains
justified. This clarification does not require editing the archived text,
but it should accompany future citations of that text.

\subsection*{Verification by the reader}
First check the precise model and its coefficient field; then run the
rectifier certificates; next check the constructive passage to the kernel
and slice. For new general statements, the principal object to verify
remains the mathematical proof. If an error is found, please identify the model,
the statement, the execution command and the full log. The authors invite
readers to undertake this verification and to develop the research programme.
